\begin{pastexam}[label={px:ntut-113-ci-02}]{北科大資工所 113 學年度 第二題}
（是非；只寫答案）

Given a set $S$ of characters ($A$, $B$, $C$, $D$, $E$, $F$, $G$, $H$, $I$) with relative frequencies (10, 5,
23, 18, 7, 12, 15, 4, 6). Let $T$ be an optimal Huffman code tree for $S$. Mark by T(=true) or F(=false) each
of the following: (10 pts)
\begin{enumerate}[label=(\arabic*)]
\item $T$ is unique. (2 pts)
\item Characters $B$ and $H$ have the same path length from the root in $T$. (2 pts)
\item Recall that the \textit{weighted path length} (WPL) is defined as $\sum(f_i\times p_i)$ where $f_i$ is
the frequency of character $i$ and $p_i$ is the length of the path in $T$ from the root to the character $i$.
Then, the weighted path length of $T$ is 249. (2 pts)
\item The path length of character $C$ is 2. (2 pts)
\item Consider another set $S'$ of characters ($A$, $B$, $C$, $D$, $E$, $F$, $G$, $H$) with relative
frequencies (10, 5, 23, 18, 13, 12, 15, 4). Let $T'$ be an optimal Huffman code tree for $S'$. Then,
WPL($T$) and WPL($T'$) differ by 13. (2 pts)
\end{enumerate}
\end{pastexam}

\begin{pastsoln}
答案：(1) F　(2) T　(3) F　(4) T　(5) T。

\textbf{建樹。} 照資結書第~19.3~節，每次取出最小的兩個合併（總頻率 $100$）：
\begin{center}
\small
\begin{tabular}{clll}
\toprule
步 & 合併 & 新節點 & 剩下的權重 \\
\midrule
1 & $H(4) + B(5)$ & 9 & 6, 7, 9, 10, 12, 15, 18, 23 \\
2 & $I(6) + E(7)$ & 13 & 9, 10, 12, 13, 15, 18, 23 \\
3 & $9 + A(10)$ & 19 & 12, 13, 15, 18, 19, 23 \\
4 & $F(12) + 13$ & 25 & 15, 18, 19, 23, 25 \\
5 & $G(15) + D(18)$ & 33 & 19, 23, 25, 33 \\
6 & $19 + C(23)$ & 42 & 25, 33, 42 \\
7 & $25 + 33$ & 58 & 42, 58 \\
8 & $42 + 58$ & 100 & 100 \\
\bottomrule
\end{tabular}
\end{center}
每一步最小的兩個都唯一，沒有平手。 由合併關係讀出每個字元的深度（根到它的路徑長）：根的兩邊是 $42$、$58$；
$42 = 19 + C$，$19 = 9 + A$，$9 = H + B$；$58 = 25 + 33$，$25 = F + 13$，$13 = I + E$，$33 = G + D$。
\begin{center}
\begin{tabular}{l*{9}{c}}
\toprule
字元 & $A$ & $B$ & $C$ & $D$ & $E$ & $F$ & $G$ & $H$ & $I$ \\
頻率 & 10 & 5 & 23 & 18 & 7 & 12 & 15 & 4 & 6 \\
深度 & 3 & 4 & 2 & 3 & 4 & 3 & 3 & 4 & 4 \\
\bottomrule
\end{tabular}
\end{center}

\textbf{(1) F。} 建樹過程雖然沒有平手（每個字元的深度因此是唯一確定的），但每個內部節點的左右子樹都可以對調，
得到不同的樹（不同的編碼），WPL 完全一樣，所以最佳的霍夫曼樹不唯一。 一般來說，頻率有平手時連各字元的深度
都可能不同（資結書第~19.3~節）。

\textbf{(2) T。} $H$、$B$ 在第 1 步合併成兄弟，深度都是 $4$。

\textbf{(3) F。}
\[
\mathrm{WPL}(T) = 10 \cdot 3 + 5 \cdot 4 + 23 \cdot 2 + 18 \cdot 3 + 7 \cdot 4 + 12 \cdot 3 + 15 \cdot 3 + 4 \cdot 4
+ 6 \cdot 4 = 299,
\]
不是 $249$。 另一種算法：WPL 等於所有內部節點的權重和（每個字元的頻率在它的每個祖先都被算一次），
$9 + 13 + 19 + 25 + 33 + 42 + 58 + 100 = 299$，一致。

\textbf{(4) T。} $C$ 是 $42$ 的子節點，$42$ 是根的子節點，深度 $2$。

\textbf{(5) T。} $S'$ 就是把 $S$ 的 $E(7)$、$I(6)$ 換成一個頻率 $13$ 的字元 $E$。 $S'$ 的霍夫曼樹：
$H + B = 9$，$9 + A = 19$，$F(12) + E(13) = 25$，$G + D = 33$，$19 + C = 42$，$25 + 33 = 58$，$42 + 58 = 100$。
內部節點權重和是 $9 + 19 + 25 + 33 + 42 + 58 + 100 = 286$。 它正好是 $T$ 拿掉內部節點 $13$（$I$、$E$ 的父節點）
的樣子，所以
\[
\mathrm{WPL}(T) - \mathrm{WPL}(T') = 299 - 286 = 13.
\]
$I$、$E$ 不是頻率最小的兩個，差剛好是 $13$ 是因為 $S'$ 的合併過程和 $S$ 一步對一步相同，只少了「$I + E$」那一步
（在 $S$ 裡 $13$ 被造出來之後，和 $S'$ 的 $E(13)$ 扮演同樣的角色）；不能直接套資結書引理~19.4，那條引理講的是
合併\textbf{最小}的兩個。
\end{pastsoln}
